\documentclass[11pt]{article}
\usepackage{ochamath}
\subjectcolor{2F5DA8}
\setsheet{線形代数・電気系院試補充}{複素ベクトルの内積　演習}{https://university-math-crj.pages.dev/linear-algebra/complex-inner-products/}
\begin{document}
\makesheethead{解答}
自作問題。本文の例題とは数値や設定が異なります。条件・途中式・検算を示してください。
\problem[共役内積]
$x=(1,2j)^{\mathsf T},y=(2j,1)^{\mathsf T}$ の内積、長さ、直交性を求めよ。
\begin{ans}
$x^*y=2j-2j=0$。$x^*x=y^*y=5$ なので両方の長さは $\sqrt5$ で直交。単なる転置なら $x^{\mathsf T}y=4j$ となって内積の規約と一致しない。成分の共役を取ってから掛ける。
\end{ans}
\point{虚数単位の共役は $-j$。}
\problem[Cauchy--Schwarz]
複素内積について $|x^*y|\le\|x\|\|y\|$ と等号条件を証明せよ。
\begin{ans}
$y\ne0$ なら $c=y^*x/(y^*y)$、$r=x-cy$ とする。$y^*r=0$、$0\le\|r\|^2=\|x\|^2-|y^*x|^2/\|y\|^2$ より不等式。等号は $r=0$ で、$x$ が $y$ の複素スカラー倍。$y=0$ では両辺0で等号。
\end{ans}
\point{複素係数の射影で証明する。}
\newpage
\problem[複素Gram--Schmidt]
$u_1=(1,-j),u_2=(0,1)$ を正規直交化せよ。
\begin{ans}
$q_1=(1,-j)/\sqrt2$、$q_1^*u_2=j/\sqrt2$。$v_2=u_2-q_1(j/\sqrt2)=(-j/2,1/2)$、長さ $1/\sqrt2$。$q_2=(-j,1)/\sqrt2$。$q_1^*q_2=(-j+j)/2=0$、両長さ1で検算できる。
\end{ans}
\point{射影係数は先頭のベクトルを共役にする。}
\problem[射影行列]
$q=(1,-j)/\sqrt2$ への射影行列と、$b=(1,0)$ の射影を求めよ。
\begin{ans}
$P=qq^*=\frac12\begin{pmatrix}1&j\\-j&1\end{pmatrix}$。$Pb=(1/2,-j/2)$、残差 $b-Pb=(1/2,j/2)$。$q^*(b-Pb)=0$、$P^2=P,P^*=P$。射影成分と残差のノルム2乗はそれぞれ1/2で合計 $\|b\|^2=1$。
\end{ans}
\point{Hermite性と冪等性を両方検算する。}
\newpage
\problem[複素最小二乗]
$a=(1,j,1)^{\mathsf T},b=(0,1,1)^{\mathsf T}$ に対し $\min_{c\in\mathbb C}\|ac-b\|^2$ を解け。
\begin{ans}
$a^*a=3,a^*b=1-j$。正規方程式 $3c=1-j$ から $c=(1-j)/3$。残差は $a^*(ac-b)=0$、最小誤差2乗は $2-|1-j|^2/3=4/3$。$a$ は非零なので1列としてフルランクで唯一解。
\end{ans}
\point{解の一意性と残差の直交も示す。}
\problem[最大値フェーザの電力]
最大値フェーザ $V=2,I=1-j$ の平均有効電力を求め、実効値規約との違いを述べよ。
\begin{ans}
最大値規約では $P=\frac12\operatorname{Re}(V\overline I)=\frac12\operatorname{Re}(2+2j)=1$。実効値フェーザは各値を $\sqrt2$ で割り、式の係数1/2を除くので同じ1になる。与えられた最大値の数値を実効値の式へそのまま入れると2となり誤り。
\end{ans}
\point{フェーザの大きさの規約を固定する。}
\end{document}
